\begin{figure}[ht]
\centering
\begin{tikzpicture}[x=1.0cm, y=1.6cm, font=\footnotesize]

% A peaked integrand on [0,1] with the recursive bisection pattern of
% adaptive Simpson drawn underneath: the interval is subdivided densely
% where the integrand is curved (near the peak) and left coarse where
% the integrand is flat. Tick depth encodes recursion depth.

% the curve: a Lorentzian-like peak centred at x=0.7
% f(x) = 1 / (1 + 400 (x-0.7)^2 ), scaled to fit
\draw[->, thick] (-0.2,0) -- (10.6,0) node[right] {$x$};
\draw[->, thick] (0,-0.05) -- (0,1.25) node[above] {$f(x)$};
\node[anchor=north, font=\scriptsize] at (0,-0.05) {$0$};
\node[anchor=north, font=\scriptsize] at (10,-0.05) {$1$};

% peak curve (x in screen units 0..10 maps to [0,1]; peak at screen 7)
\draw[thick, engblue] plot[domain=0:10, samples=120] (\x,
  {1/(1 + 4*(\x-7)*(\x-7))});

% Coarse subdivisions away from peak (depth 1): big intervals
\foreach \xb in {0,2.5} {
  \draw[engred, thick] (\xb,0) -- (\xb,-0.14);
}
% Medium subdivisions approaching peak (depth 2)
\foreach \xb in {5.0,6.0} {
  \draw[engred, thick] (\xb,0) -- (\xb,-0.20);
}
% Fine subdivisions at the peak (depth 3..4)
\foreach \xb in {6.5,6.75,7.0,7.25,7.5} {
  \draw[engred, thick] (\xb,0) -- (\xb,-0.26);
}
% Medium again on the far flank (depth 2)
\foreach \xb in {8.0,9.0} {
  \draw[engred, thick] (\xb,0) -- (\xb,-0.20);
}
% Endpoint
\draw[engred, thick] (10,0) -- (10,-0.14);

\node[engred, font=\scriptsize, anchor=north] at (1.25,-0.34)
  {coarse: flat region};
\node[engred, font=\scriptsize, anchor=north] at (7.0,-0.40)
  {fine: high curvature};

\end{tikzpicture}
\caption[Adaptive-quadrature subdivision following the integrand]{An
adaptive Simpson routine on a peaked integrand. The routine estimates the
local error on each panel by comparing a Simpson estimate to the sum of two
Simpson estimates on the halves, and recurses only where the discrepancy
exceeds the per-panel tolerance. The tick depth marks the recursion depth:
the flat flanks are accepted at coarse resolution after one or two
subdivisions, while the curved peak draws repeated bisection until the local
error falls below tolerance. The total work concentrates where the integrand
has structure, which is the property fixed-order rules lack. The worked trace
in the text carries one such subdivision numerically. Source: authors'
analysis.}
\label{fig:vol02:ch16:adaptive-quadrature}
\end{figure}
